\documentclass{ximera}

% MATH -----------------------------------------------------------
\newcommand{\nats}{\mathbb N}
\newcommand{\ints}{\mathbb Z}
\newcommand{\rats}{\mathbb Q}
\newcommand{\reals}{\mathbb R}
\newcommand{\complex}{\mathbb C}
\newcommand{\powerset}{\mathscr P}

%-------------------------------------------------------------

\pgfplotsset{compat=1.18}
\begin{document}
\begin{abstract}
 \end{abstract}

    \title{Exemplar Examples 10}
    
    \maketitle

Let's consider the second order linear homogeneous differential equation $$\displaystyle y\,^{\prime\prime}+\dfrac{3}{x}y\,^{\prime}+\dfrac{1}{x^2}y=0.$$

\begin{enumerate}

\item Suppose $y_1$ and $y_2$ are a fundamental set of solutions to this differential equation on $(0,\infty)$ and that their Wronskian $W(x)$ satisfies $W(1)=1$.  Let's determine $W(x)$.

\vfill


% W=Ce^(-3ln(x))=C/x^3

\item Let's show that $y_1=x^{-1}\ln{(x)}$ is a solution to this differential equation.

% PLUG IN TO LEFT SIDE AND SIMPLIFY TO 0

\vfill

\item It turns out that $y_2=x^{-1}\ln{(x^2)}$ is also a solution.  Does $\{y_1,y_2\}$ form a fundamental set of solutions to $\displaystyle y\,^{\prime\prime}+\dfrac{3}{x}y\,^{\prime}+\dfrac{1}{x^2}y=0$ on $(0,\infty)$?

\vfill

% No, because y_2 = 2y_1

\newpage

\item It turns out that with $y_2=x^{-1}$ we get a fundamental set of solutions. Let's solve the IVP $\displaystyle y^{\,\prime\prime}+\dfrac{3}{x}y^{\,\prime}+\dfrac{1}{x^2}y=0$, $y(1)=1$, $y^{\,\prime}(1)=0$.


% y=x^{-1}ln(x)+x^{-1}


\end{enumerate}

\end{document}